\begin{trapbox}

	\begin{description}[style=nextline, leftmargin=2.8em, labelsep=0.5em, itemsep=0.35em]
		\item[\textbf{\textcolor{CTrap}{易错点一（通径遗漏）：}}]
			在证明或应用该结论时，只考虑斜率存在的焦点弦，未验证通径（即$AB\perp x$轴）情况，导致推导不完整。
		\item[\textbf{\textcolor{CTrap}{正确理解（分类讨论）：}}]
			必须单独验证通径情形，或采用$x=my+\frac{p}{2}$形式的直线方程，自动包含斜率不存在的情况。
		\item[\textbf{\textcolor{CTrap}{错因分析（思维定势）：}}]
			习惯性设直线为$y=k(x-\frac{p}{2})$，忽略了$k$可能不存在，从而遗漏特殊情况。
	\end{description}

	\medskip
	\begin{description}[style=nextline, leftmargin=2.8em, labelsep=0.5em, itemsep=0.35em]
		\item[\textbf{\textcolor{CTrap}{易错点二（公式混用）：}}]
			计算$|AF|$时错误地使用$|AF|=x_1-\frac{p}{2}$或$|AF|=x_1$，导致结果错误。
		\item[\textbf{\textcolor{CTrap}{正确理解（定义法）：}}]
			由抛物线定义，$|AF|$等于点$A$到准线$x=-\frac{p}{2}$的距离，即$|AF|=x_1+\frac{p}{2}$。
		\item[\textbf{\textcolor{CTrap}{错因分析（记忆混淆）：}}]
			容易将椭圆的焦半径公式与抛物线混淆，或对抛物线准线位置不清楚。
	\end{description}

	\medskip
	\begin{description}[style=nextline, leftmargin=2.8em, labelsep=0.5em, itemsep=0.35em]
		\item[\textbf{\textcolor{CTrap}{易错点三（符号忽略）：}}]
			盲目套用结论，忽略$p>0$的前提条件；当抛物线开口向左（$y^2=-2px$）时仍代入$\frac{2}{p}$形式。
		\item[\textbf{\textcolor{CTrap}{正确理解（条件限制）：}}]
			结论仅适用于标准方程$y^2=2px\,(p>0)$，此时焦点在$x$轴正半轴；若开口方向不同，结论形式会变化。
		\item[\textbf{\textcolor{CTrap}{错因分析（忽视前提）：}}]
			未注意抛物线方程中$p$的符号意义，死记结论，不关注适用条件。
	\end{description}

\end{trapbox}
